常识验证要求模型判断一个自然语言陈述是否符合现实世界知识，是评估语言模型理解能力的重要任务。本文以 SemEval-2020 Task 4 \comve 子任务 A 为实验平台，研究大语言模型在两句相近文本中识别非常识句子的可靠性。除 RoBERTa 微调、零样本提示、少样本提示和自一致性等基础方法外，我们从两个互补方向扩展分析：一是评估模型在句子顺序、提示模板、少样本示例和采样随机性变化下的鲁棒性与一致性；二是探索 constraint-first prompting、generate-verify-rerank 和 consistency-based uncertainty 等轻量推理时增强方法。Task A 结果显示，RoBERTa 微调达到 0.868 准确率和 0.939 顺序一致性；最佳 LLM 为 Qwen3-1.7B thinking zero-shot direct prompt，准确率为 0.952，顺序一致性为 0.929，提示一致性为 0.849。8-shot 示例没有带来整体收益；no-thinking self-consistency 能将 Qwen3.5-2B minimal prompt 从 0.684 提升到 0.720，但仍低于 thinking zero-shot。Task B 结果显示，验证增强策略在 stronger models 上有进一步增益，例如 Qwen3.5-2B thinking 从 direct 的 0.782 提升到 rerank 的 0.818。整体来看，标准准确率不足以完整刻画大语言模型的常识能力；更可靠的常识验证还需要同时关注回答稳定性、解释是否支持判断，以及模型能否识别和利用自身不确定性。
